\section{预训练目标构造}
\label{sec:nlp-llm-objectives}

概率计算本身没有规定应当学会什么。相同网络可以只预测下一词元，也可以根据缺口两侧恢复中间片段；相同文本也可以产生多种监督记录。构造目标时，应分别回答：预测时允许看到哪些信息，要预测哪个对象，以及哪些错误按多大权重计入损失。三项决定需要共同明确，才能判断一个结构是否泄露答案、一个损失是否对应使用时的问题。

\subsection{可见条件与信息范围}
\label{sec:nlp-llm-objective-visibility}

对完整教学序列“读书会／改到／周五／下午／两点／举行”，若要学习左到右续写，预测“周五”时只能看到前两个词元。若要学习填空，可以把“周五”遮住，并同时提供左右两侧。两者预测对象看似相同，条件分布却不同：$p_\theta(x_3\mid x_1,x_2)$与$p_\theta(x_3\mid x_1,x_2,x_4,x_5,x_6)$不是同一个量。条件增加后问题可能更容易，但其损失也不能直接作为前一问题的替代。

设$M_{ij}\in\{0,1\}$表示输入位置$i$能否读取位置$j$。在注意力分数中，对$M_{ij}=0$的位置加上$-\infty$，再对允许位置归一化，就实现了可见范围。因果掩码取$M_{ij}=1$当且仅当$j\leq i$，位置$i$随后预测$x_{i+1}$；前缀方案可以让一段已给条件相互读取，同时限制待生成部分只能读取条件和自身过去。掩码的作用是限制信息流，而不是标记哪些位置计分，两者要分别设置。

另一种改变条件的方法是重排文本。例如先写出前缀与后缀，最后才放被移走的中间片段，那么普通因果计算在预测中间片段时，已经能够读取它两侧的内容。这里没有让网络读取计算序列中的未来，而是改变了原文位置与计算顺序的对应关系。训练记录必须保留特殊边界标记，否则模型无法判断当前是在续写原文还是补写缺口。

泄漏还可能发生在多层传播中。即使某个输出位置没有直接读取目标位置，若它先读取一个已经包含目标信息的中间表示，仍可能间接得到答案。因此，信息范围应针对整个计算图检查，不能只观察最后一层的掩码。需要排除的是目标通过本应不可见的位置或计算路径泄漏，而不是禁止一切复制。
例如从可见材料提取实体，或预测前文已经出现过的词元，都可以是合法的条件预测；
关键在于预测时使用的信息是否符合当前目标约定。

\subsection{预测对象与监督位置}
\label{sec:nlp-llm-objective-positions}

在可见条件确定后，还需要说明输出对应哪些目标。下一词元预测把输入位置$t$的表示对应到目标$x_{t+1}$；被遮蔽片段恢复把缺失位置对应到原始内容；多个未来位置预测则让同一前缀产生不同步长的输出。目标的编号必须与可见范围一致，不能把输入$x_t$对应到不受遮蔽的目标$x_t$，再误以为得到了语言预测任务。

令$\mathcal S(x)$表示一个样本中参与计分的目标位置集合，$C_i(x)$表示预测目标$x_i$时允许使用的条件。样本组织与目标构造共同产生记录$\{(C_i(x),x_i):i\in\mathcal S(x)\}$。在下一词元目标中，可以令$\mathcal S(x)=\{1,\ldots,T\}$、$C_i(x)=x_{<i}$；在遮蔽目标中，$\mathcal S(x)$通常只包含被遮蔽位置，条件包括未遮蔽内容及其位置。开始、结束、填充等特殊符号是否计分，也需要在这个集合中明确。

例如，把两篇文档装入一个样本后，填充位置通常不应作为普通正文目标计分；若文档边界强制隔离，第二篇文档的第一个词元也不能依赖第一篇的尾部。对重叠切分的长文档，还可以只对新引入部分计分，让重叠前缀提供条件而不重复贡献全部损失。改变计分位置会改变每段材料的实际学习权重，即使输入字符串没有变化。

多个未来位置预测还要区分辅助目标与生成协议。让当前位置同时预测$x_{t+1}$和$x_{t+2}$，会使共享表示收到两项误差信号；但若两个输出头都只读取$x_{1:t}$，第二个头就没有以第一个头实际选出的词元为条件。它们是给定同一前缀的两个预测分布，不能仅因共同训练就称为完整的未来联合分布。相应结构将在\ref{sec:nlp-llm-objective-comparison}小节写出。

\subsection{概率损失与目标组合}
\label{sec:nlp-llm-objective-loss}

现在可以用目标词元得到的概率为每项预测计分。若一次预测给真实目标分配的概率为$p$，负对数损失$-\log p$在$p$接近1时较小，在$p$接近0时增大。对一个样本及其可能的随机遮蔽，通用的计分形式为
\begin{equation}
 \ell(\theta;x)=
 -\frac{1}{Z(x)}\sum_{i\in\mathcal S(x)}
 w_i(x)\log p_\theta(x_i\mid C_i(x)),
 \label{eq:nlp-llm-objective-general}
\end{equation}
其中$w_i(x)\geq0$为位置权重，$Z(x)>0$给出归一化约定。若所有目标位置等权且$Z(x)=|\mathcal S(x)|$，得到每样本的平均目标词元损失；若先汇总一批样本的所有目标词元再平均，长样本对批次目标的贡献更大。两种口径需要明确，不能把它们混称为相同的平均损失。

对实际采样分布$P_{\bm\alpha}$及目标构造中的随机性取期望，得到学习的总体目标$\mathcal L(\theta)=\mathbb E[\ell(\theta;x)]$。训练时以抽到的有限批次估计其梯度。对于条件分布$q(\cdot\mid C)$，等权交叉熵可分解为该分布自身的熵与$D_{\mathrm{KL}}(q\|p_\theta)$之和，因此在固定条件、充分表达且能够求得最优值时，目标鼓励模型拟合数据条件分布。现实中$q$未知，有限材料与有限优化使这个解释不能变成对任意新输入的保证。

若同时使用若干目标，需说明每一项计算什么，再设置组合系数。例如$\mathcal L=\mathcal L_{\mathrm{next}}+\lambda\mathcal L_{\mathrm{future}}$让未来预测作为辅助信号，$\lambda\geq0$决定其相对作用。改变$\lambda$不仅改变数值大小，也改变共享参数收到的更新方向；一项损失下降可能伴随另一项上升。只比较总目标值，还可能掩盖各部分之间的取舍。

损失值还受到词元化、可见条件、目标位置和归一化的影响。具有双侧条件的填空模型得到较低损失，不足以说明它在左到右生成时更准确。即使协议完全相同，平均预测损失也可能被常见局部搭配主导，而用户关心的是少数关键条件或精确推理。预训练目标提供可学习的代理，具体能力需要独立检查。

\subsection{代表目标的完整对照}
\label{sec:nlp-llm-objective-comparison}

先固定同一序列$x_{1:6}$为上述通知短句。下一词元预测（next-token prediction, NTP）对各前缀预测下一个目标，损失可取$-\sum_{i=1}^{6}\log p_\theta(x_i\mid x_{<i})$，并按计分词元数平均。预测“周五”时看不到“下午两点”，预测“举行”时可以看到此前全部内容。所有目标共享一套概率映射，数据提供了大量不同长度的前缀。

中间填充（fill-in-the-middle, FIM）把原文分成前缀$A$、中间$B$和后缀$C$。若取$A=$“读书会改到”、$B=$“周五”、$C=$“下午两点举行”，可构造带边界标记的计算序列“前缀标记、$A$、后缀标记、$C$、中间标记、$B$”。在该序列上继续使用下一词元损失，预测$B$时就能读取$A$与$C$。FIM改变的是样本顺序与条件组织，仍可沿用因果结构；其混入比例与切分方式会影响学习结果\scite{nlp-llm-fim2022}。

UL2把多种去噪记录混合：常规跨度遮蔽只移走部分片段，顺序模式把后段作为待恢复内容，强遮蔽模式则隐藏更多内容或更长片段。它在编码器—解码器结构中，用模式标记说明当前恢复要求，再预测被移走内容\scite{nlp-llm-ul2-2022}。例如，输入“读书会改到、缺口、举行”，目标可为带片段标记的“周五下午两点”；换成顺序模式时，输入只保留前缀。这里改变了损坏方式与可见范围，不能把UL2理解为给普通NTP损失换了一个系数。

\term{多词元预测}（multi-token prediction, MTP）让同一前缀还承担更远未来的预测。共享主干、独立未来头的一种目标为
\begin{equation}
 \mathcal L_{\mathrm{MTP}}=
 -\sum_{j=1}^{J}\lambda_j
 \sum_{t:\,t+j\leq T}
 \log p_{\theta,j}(x_{t+j}\mid x_{1:t}),
 \label{eq:nlp-llm-mtp}
\end{equation}
其中$j$是预测步长，$J$是头数，$\lambda_j\geq0$为权重；这里省略了样本平均。以$x_{1:2}=$“读书会改到”为条件，两个头分别预测“周五”和“下午”。后一个头必须从同一前缀承担更远预测，因此对主干表示提出了额外要求，但两个头的边缘概率相乘并不自动给出真实未来联合分布\scite{nlp-llm-mtp2024}。

DeepSeek-V3采用顺序连接的MTP模块，将前一模块表示与后续词元嵌入共同用于更远预测，不应与独立未来头的条件完全混同\scite{deepseek2024v3}。辅助未来目标在部分代码、生成及算法任务上提供了收益证据，但收益受到模型规模、训练量和任务类型影响，并非增加预测头就必然改善所有任务。使用这些头提出多个候选词元，还需要\ref{chap:nlp-model-inference}章的验证与接受规则，训练目标本身不决定生成分布是否保持不变。

掩码扩散语言模型（masked diffusion language model）采用另一条生成路线：先规定怎样把完整文本逐步变成包含掩码的序列，再学习反向恢复。为说明一种常见设定，取噪声水平$\tau\sim\operatorname{Uniform}(0,1)$，将每个词元独立地以概率$\tau$替换为掩码符号，得到$\tilde x$。在线性遮蔽日程下，可用如下加权恢复目标：
\begin{equation}
 \mathcal L_{\mathrm{mask}}=
 \mathbb E_{\tau,x,\tilde x}
 \left[-\frac{1}{\tau}
 \sum_{i:\,\tilde x_i=\mathrm{mask}}
 \log p_\theta(x_i\mid\tilde x)\right].
 \label{eq:nlp-llm-masked-diffusion}
\end{equation}
期望覆盖完整文本、噪声水平及随机掩码。权重与所选前向损坏日程有关，不能任意移植到所有扩散模型；若使用不同日程、离散步数或预测参数化，就需要相应调整目标。

MDLM研究了掩码扩散的简化参数化与恢复损失，LLaDA在较大语言模型中采用掩码恢复路线，Dream则进一步探索扩散语言模型的学习与生成\scite{nlp-llm-mdlm2024,nlp-llm-llada2025,nlp-llm-dream2025}。对通知例子，受损输入可以同时保留“周五”和“举行”，要求恢复中间时间；当接近全掩码时，模型则必须更多依赖已学的文本分布。生成时通常从全掩码或部分已知序列出发，反复恢复与更新，网络输出的恢复分布并非式\eqref{eq:nlp-llm-sequence-probability}中的前缀条件分布，因此不能直接把这些输出按左到右顺序连乘来计算文本概率。完整采样过程留到\ref{chap:nlp-model-inference}章，当前的结论是：损坏过程、恢复目标和生成过程必须配套定义。
